% Options for packages loaded elsewhere
\PassOptionsToPackage{unicode}{hyperref}
\PassOptionsToPackage{hyphens}{url}
\PassOptionsToPackage{dvipsnames,svgnames,x11names}{xcolor}
%
\documentclass[
  11pt,
]{article}
\usepackage{amsmath,amssymb}
\usepackage{iftex}
\ifPDFTeX
  \usepackage[T1]{fontenc}
  \usepackage[utf8]{inputenc}
  \usepackage{textcomp} % provide euro and other symbols
\else % if luatex or xetex
  \usepackage{unicode-math} % this also loads fontspec
  \defaultfontfeatures{Scale=MatchLowercase}
  \defaultfontfeatures[\rmfamily]{Ligatures=TeX,Scale=1}
\fi
\usepackage{lmodern}
\ifPDFTeX\else
  % xetex/luatex font selection
\fi
% Use upquote if available, for straight quotes in verbatim environments
\IfFileExists{upquote.sty}{\usepackage{upquote}}{}
\IfFileExists{microtype.sty}{% use microtype if available
  \usepackage[]{microtype}
  \UseMicrotypeSet[protrusion]{basicmath} % disable protrusion for tt fonts
}{}
\makeatletter
\@ifundefined{KOMAClassName}{% if non-KOMA class
  \IfFileExists{parskip.sty}{%
    \usepackage{parskip}
  }{% else
    \setlength{\parindent}{0pt}
    \setlength{\parskip}{6pt plus 2pt minus 1pt}}
}{% if KOMA class
  \KOMAoptions{parskip=half}}
\makeatother
\usepackage{xcolor}
\usepackage[margin=22mm]{geometry}
\usepackage{listings}
\newcommand{\passthrough}[1]{#1}
\lstset{defaultdialect=[5.3]Lua}
\lstset{defaultdialect=[x86masm]Assembler}
\usepackage{longtable,booktabs,array}
\usepackage{calc} % for calculating minipage widths
% Correct order of tables after \paragraph or \subparagraph
\usepackage{etoolbox}
\makeatletter
\patchcmd\longtable{\par}{\if@noskipsec\mbox{}\fi\par}{}{}
\makeatother
% Allow footnotes in longtable head/foot
\IfFileExists{footnotehyper.sty}{\usepackage{footnotehyper}}{\usepackage{footnote}}
\makesavenoteenv{longtable}
\setlength{\emergencystretch}{3em} % prevent overfull lines
\providecommand{\tightlist}{%
  \setlength{\itemsep}{0pt}\setlength{\parskip}{0pt}}
\setcounter{secnumdepth}{5}
\usepackage{xcolor}
\usepackage{fancyhdr}
\usepackage{titlesec}
\definecolor{AcademicBlue}{HTML}{17365D}
\definecolor{CodeBackground}{HTML}{F4F6F8}
\pagestyle{fancy}
\fancyhf{}
\fancyhead[L]{\small COMP3028 -- Software Construction}
\fancyhead[R]{\small Case Study 2}
\fancyfoot[C]{\thepage}
\setlength{\headheight}{14pt}
\titleformat{\section}{\large\bfseries\color{AcademicBlue}}{\thesection.}{0.5em}{}
\lstdefinelanguage{JavaScript}{keywords={import,from,const,let,export,async,await,function,return,try,catch,throw,new},sensitive=true,morecomment=[l]{//},morestring=[b]',morestring=[b]",morestring=[b]`}
\lstset{basicstyle=\ttfamily\footnotesize,breaklines=true,columns=fullflexible,keepspaces=true,showstringspaces=false,backgroundcolor=\color{CodeBackground},frame=single,rulecolor=\color{black!15},keywordstyle=\bfseries\color{AcademicBlue}}
\ifLuaTeX
  \usepackage{selnolig}  % disable illegal ligatures
\fi
\IfFileExists{bookmark.sty}{\usepackage{bookmark}}{\usepackage{hyperref}}
\IfFileExists{xurl.sty}{\usepackage{xurl}}{} % add URL line breaks if available
\urlstyle{same}
\hypersetup{
  pdftitle={Code Review and Refactoring},
  pdfauthor={COMP3028 --- Software Construction},
  colorlinks=true,
  linkcolor={Maroon},
  filecolor={Maroon},
  citecolor={Blue},
  urlcolor={Blue},
  pdfcreator={LaTeX via pandoc}}

\title{Code Review and Refactoring}
\usepackage{etoolbox}
\makeatletter
\providecommand{\subtitle}[1]{% add subtitle to \maketitle
  \apptocmd{\@title}{\par {\large #1 \par}}{}{}
}
\makeatother
\subtitle{Case Study 2: Laboratory Equipment Loan Service}
\author{COMP3028 --- Software Construction}
\date{}

\begin{document}
\maketitle

\textbf{Assessment mode:} Individual implementation with paired peer
review\\
\textbf{Total marks:} 20\\
\textbf{Submission:} Source code, tests, pull-request evidence, and a
concise review report\\
\textbf{Deadline and target branch:} As announced by the instructor

\hypertarget{context}{%
\section{Context}\label{context}}

A university laboratory lends equipment, such as digital multimeters and
oscilloscopes, to eligible students. A prototype service records a loan
and calculates the date on which the equipment must be returned.

The prototype was developed quickly. Although a typical seven-day
request appears to work, it does not consistently enforce lending rules
or communicate failures. You have been assigned to review and improve
the service. Your changes must undergo peer review before integration
into the designated development branch.

This exercise uses fictional records and an in-memory repository. No
database, credentials, network access, or external packages are
required. Changes are not retained after the process exits.

\hypertarget{intended-learning-outcomes}{%
\section{Intended learning outcomes}\label{intended-learning-outcomes}}

On completion, students should be able to:

\begin{enumerate}
\def\labelenumi{\arabic{enumi}.}
\tightlist
\item
  Identify defects and maintainability concerns in asynchronous
  JavaScript.
\item
  Implement validation and enforce explicit domain rules.
\item
  Calculate deterministic dates without mutating caller-owned objects.
\item
  Preserve successful outcomes while documenting intentional behaviour
  changes.
\item
  Test success, failure, and boundary conditions using mocks.
\item
  Provide constructive, evidence-based feedback through a pull request.
\end{enumerate}

\hypertarget{supplied-implementation}{%
\section{Supplied implementation}\label{supplied-implementation}}

The primary review target is
\passthrough{\lstinline!src/loanService.js!}. The following listing is
the supplied implementation.

\begin{lstlisting}
// Intentionally incomplete student starter. See docs/assignment.md.
// Repository and clock dependencies are provided by the caller.
export async function borrowEquipment(request, repository, clock = () => new Date()) {
    const student = await repository.findStudent(request.studentId);
    const equipment = await repository.findEquipment(request.equipmentId);

    const dueDate = clock();
    dueDate.setUTCDate(dueDate.getUTCDate() + 7);

    try {
        const loan = await repository.createLoan({
            studentId: student.id,
            equipmentId: equipment.id,
            dueAt: dueDate.toISOString()
        });
        return loan;
    } catch (error) {
        return null;
    }
}
\end{lstlisting}

The implementation already awaits repository operations, returns a
stored record on success, and accepts substitutable dependencies.
Reviewers must acknowledge these strengths rather than assuming that
every part of the code is defective.

\hypertarget{interface-and-repository-contract}{%
\section{Interface and repository
contract}\label{interface-and-repository-contract}}

Preserve the public interface:

\begin{lstlisting}
borrowEquipment(request, repository, clock)
\end{lstlisting}

The request contains \passthrough{\lstinline!studentId!},
\passthrough{\lstinline!equipmentId!}, and
\passthrough{\lstinline!days!}. The optional clock returns a valid Date
and defaults to the current time. It is trusted infrastructure; students
need not validate arbitrary clock implementations.

\begin{longtable}[]{@{}
  >{\raggedright\arraybackslash}p{(\columnwidth - 2\tabcolsep) * \real{0.5000}}
  >{\raggedright\arraybackslash}p{(\columnwidth - 2\tabcolsep) * \real{0.5000}}@{}}
\toprule\noalign{}
\begin{minipage}[b]{\linewidth}\raggedright
Repository method
\end{minipage} & \begin{minipage}[b]{\linewidth}\raggedright
Contract
\end{minipage} \\
\midrule\noalign{}
\endhead
\bottomrule\noalign{}
\endlastfoot
findStudent(id) & Resolves to a record containing id and active, or null
if absent. \\
findEquipment(id) & Resolves to a record containing id and available, or
null if absent. \\
createLoan(data) & Resolves to the stored loan, including its generated
id; rejects on storage failure or unavailable equipment. \\
\end{longtable}

Creation receives exactly the business fields
\passthrough{\lstinline!studentId!},
\passthrough{\lstinline!equipmentId!}, and
\passthrough{\lstinline!dueAt!}. The repository owns the atomic
operation that checks availability, creates the loan, and marks the
equipment unavailable. The service must not implement a separate
availability update.

The service must check eligibility before attempting creation. A
repository may still reject creation if availability changes
subsequently; such an error must reach the caller.

\hypertarget{task-a-initial-code-review}{%
\section{Task A: Initial code review}\label{task-a-initial-code-review}}

Identify at least five concerns. For each concern, report the relevant
location, the observed issue, its possible consequence, and a justified
improvement.

Consider input types, missing records, eligibility, requested duration,
ownership of mutable Date objects, and the treatment of repository
failures. Distinguish an actual defect from a stylistic preference.

Record your findings before implementation. Do not limit your review to
the failures already exposed by the provided tests.

\hypertarget{task-b-required-behaviour}{%
\section{Task B: Required behaviour}\label{task-b-required-behaviour}}

Apply validation and eligibility checks in the order specified below.
Where more than one condition fails, report the first applicable error.

\begin{longtable}[]{@{}
  >{\raggedright\arraybackslash}p{(\columnwidth - 4\tabcolsep) * \real{0.3333}}
  >{\raggedright\arraybackslash}p{(\columnwidth - 4\tabcolsep) * \real{0.3333}}
  >{\raggedright\arraybackslash}p{(\columnwidth - 4\tabcolsep) * \real{0.3333}}@{}}
\toprule\noalign{}
\begin{minipage}[b]{\linewidth}\raggedright
Order
\end{minipage} & \begin{minipage}[b]{\linewidth}\raggedright
Condition
\end{minipage} & \begin{minipage}[b]{\linewidth}\raggedright
Required result
\end{minipage} \\
\midrule\noalign{}
\endhead
\bottomrule\noalign{}
\endlastfoot
1 & Request is missing, null, an array, or not an object & Reject with
INVALID\_INPUT; no repository access. \\
1 & Either ID is not a positive JavaScript safe integer, or days is not
an integer from 1 to 14 inclusive & Reject with INVALID\_INPUT; no
repository access. \\
2 & Student is absent & Reject with STUDENT\_NOT\_FOUND. \\
3 & Student is not active & Reject with STUDENT\_INACTIVE. \\
4 & Equipment is absent & Reject with EQUIPMENT\_NOT\_FOUND. \\
5 & Equipment is unavailable & Reject with EQUIPMENT\_UNAVAILABLE. \\
6 & All conditions pass & Calculate the due date, await one creation
operation, and return the stored loan. \\
\end{longtable}

For valid records, active and available are Boolean fields. Numeric
strings such as ``1'' must be rejected rather than converted. Ignore
additional request fields; do not pass them to storage.

Each specified domain or validation error must be an Error instance with
a \passthrough{\lstinline!code!} property matching the table and a
meaningful message. Exact wording is not prescribed.

Additional requirements:

\begin{itemize}
\tightlist
\item
  Do not attempt loan creation after any failed validation or
  eligibility check.
\item
  Read the clock once for a successful eligibility path. Calculate dueAt
  as exactly the requested number of 24-hour days after that instant,
  represented as an ISO UTC string.
\item
  Do not mutate the request, repository records, or Date returned by the
  clock.
\item
  Propagate failures from any repository method. Rethrow the original
  error or preserve it as the cause of a new error. Do not replace
  failure with null or a successful result.
\item
  Preserve the successful return contract: the stored loan is returned
  only after creation completes.
\item
  Preserve the public signature and the provided repository contract.
  Internal helper functions may be introduced.
\item
  Do not add database setup, external services, or unrelated
  functionality.
\end{itemize}

\hypertarget{task-c-refactoring-rationale}{%
\section{Task C: Refactoring
rationale}\label{task-c-refactoring-rationale}}

Explain the purpose and behavioural impact of the changes.

Separating validation into a helper or improving names may preserve
behaviour. Rejecting invalid requests, enforcing eligibility, respecting
requested duration, and exposing previously swallowed errors
intentionally change behaviour. Identify these changes explicitly.

Discuss why mutation of a caller-owned Date can cause surprising
behaviour in repeated calls. Explain why dependency substitution
supports deterministic tests and independent code review.

\hypertarget{task-d-testing}{%
\section{Task D: Testing}\label{task-d-testing}}

Run the supplied tests and add coverage for requirements not yet
exercised. Do not delete, skip, or weaken supplied assertions to obtain
a passing result.

The unmodified starter is expected to produce \textbf{2 passing smoke
tests and 8 failing acceptance tests}. The final submission should pass
all supplied tests and the additional tests you write.

At minimum, extend coverage to include:

\begin{enumerate}
\def\labelenumi{\arabic{enumi}.}
\tightlist
\item
  Duration boundaries 1 and 14, values outside the range, fractional
  values, and numeric strings.
\item
  Invalid IDs, including unsafe integers, and invalid request shapes;
  verify no repository access.
\item
  Missing equipment and cases where multiple eligibility checks fail,
  demonstrating error precedence.
\item
  Failures in student lookup and equipment lookup, in addition to
  creation failure.
\item
  An exact due date across a month or year boundary and preservation of
  the input Date.
\item
  Exactly one creation for an accepted request, with only the required
  fields.
\item
  A deferred creation promise, demonstrating that success is not
  returned prematurely.
\item
  Repeated calls with a deterministic clock and unchanged caller-owned
  inputs.
\end{enumerate}

Use fresh fixtures for independent tests. Assertions should verify
relevant behaviour rather than private helper names or formatting. The
supplied suite is a starting point, not a complete specification.

\hypertarget{task-e-peer-review}{%
\section{Task E: Peer review}\label{task-e-peer-review}}

Each student must complete one author role and one reviewer role.

\begin{enumerate}
\def\labelenumi{\arabic{enumi}.}
\tightlist
\item
  Create a dedicated branch and complete a self-review.
\item
  Commit logical changes with meaningful messages.
\item
  Open a pull request against the instructor-designated target branch.
  Use the provided template.
\item
  Review a peer's code, rationale, and tests against this brief.
\item
  Label feedback as a required change, suggestion, or question. Explain
  the observation and consequence; propose an action where appropriate.
\item
  Respond to substantive feedback through revision or a reasoned
  explanation.
\item
  Request re-review and record the reviewer's final decision. Follow the
  instructor's merge rules.
\end{enumerate}

The reviewer should verify correctness, readability, error propagation,
date ownership, and test quality. Avoid inventing defects or imposing
undocumented personal style preferences.

An example of constructive feedback is:

\begin{quote}
The due-date variable refers to the Date returned by the clock. Updating
it also changes the caller's object, which can shift later calculations.
Please avoid mutating that object and add a test confirming that its
timestamp remains unchanged.
\end{quote}

\hypertarget{deliverables}{%
\section{Deliverables}\label{deliverables}}

Submit:

\begin{enumerate}
\def\labelenumi{\arabic{enumi}.}
\tightlist
\item
  The improved service and any justified internal helpers.
\item
  Supplied and additional tests, with execution instructions and actual
  results.
\item
  An initial review and implementation rationale.
\item
  A link to the authored pull request and evidence of the peer review
  performed.
\item
  A completed review record and a brief reflection on feedback received.
\end{enumerate}

Use the provided review-record template where helpful. Follow the
course's academic integrity policy and attribute external material.
Report only verification that you actually performed.

\hypertarget{assessment-rubric}{%
\section{Assessment rubric}\label{assessment-rubric}}

\begin{longtable}[]{@{}
  >{\raggedright\arraybackslash}p{(\columnwidth - 2\tabcolsep) * \real{0.4286}}
  >{\raggedleft\arraybackslash}p{(\columnwidth - 2\tabcolsep) * \real{0.5714}}@{}}
\toprule\noalign{}
\begin{minipage}[b]{\linewidth}\raggedright
Criterion
\end{minipage} & \begin{minipage}[b]{\linewidth}\raggedleft
Marks
\end{minipage} \\
\midrule\noalign{}
\endhead
\bottomrule\noalign{}
\endlastfoot
Accurate initial review identifying strengths and defects with
consequences & 4 \\
Correct validation, eligibility, duration calculation, and failure
propagation & 6 \\
Maintainable structure, clear names, and preservation of interfaces and
ownership & 3 \\
Meaningful additional tests covering boundaries and asynchronous failure
behaviour & 4 \\
Constructive peer review, justified responses, and documented revisions
& 3 \\
\textbf{Total} & \textbf{20} \\
\end{longtable}

\hypertarget{scope-limitations}{%
\section{Scope limitations}\label{scope-limitations}}

Authentication, authorisation, fees, notifications, renewals, student
loan quotas, and production database transactions are outside the
assessment. The repository contract provides atomic creation;
implementing production concurrency control is not required.

\textbf{Reflection question:} Why can a successful demonstration with
one valid request provide insufficient evidence that a service is ready
to merge?

\end{document}
